\documentclass{article}
\input{boilerplate}

\begin{document}

\newcommand{\lectureTitle}{Project Proposal: Generalizable RL for Neighborhood Battery Control in CityLearn}
\newcommand{\lectureDate}{Due: March 20, 2026}

\textsc{COS435 / ECE433: Introduction to RL} \hfill \lectureDate
\vspace{1em}

\maketitle

\paragraph{Authors.} Erik Dyer (\texttt{erikdyer@princeton.edu}), Ammaar Alam (\texttt{ah0952@princeton.edu}), Grace Sun (\texttt{gs0625@princeton.edu})


\paragraph{Project type:}
3. Competition-based project (The CityLearn Challenge 2023).

\paragraph{Description}

We plan to study CityLearn, an open-source gynasium environment where agents control neighborhood batteries for homes with solar generation and varying electricity demand. During each step, the controller outputs a continuous action $a_{t}\in[-1,1]$ per building, where -1 is full discharge and +1 is full charge, to improve neighborhood-level metrics like electricity cost, carbon emissions, and load.

\begin{wrapfigure}{r}{0.5\textwidth}
    \includegraphics[width=0.5\textwidth]{figures/citylearn_diagram.png}
\end{wrapfigure}

This challenge is preferred since it is a clear sequential decision-making problem coupled with delayed effects, such as how charging now affects future choices, so the value of an action depends on how it affects later states and opportunities for control \citep{vazquez2020citylearn, citylearn2023}.

We will use the CityLearn environment and starter code as an initial framework \citep{citylearnrepo}. Our initial milestone leverage Stable-Baselines3 to implement working baseline controllers, such as CityLearn's built-in Rule-Based Controller (RBC) and RL agents like PPO and SAC. Then, we will progress to testing select changes to see if they improve performance and generalization (e.g., short-horizon forecast, decentralized and centralized control, or modified  reward design). Our primary research question is whether a relatively simple RL pipeline can outperform classical hierarchical optimization (e.g., CHESCA, 2023 winner) and rule-based strategies while remaining robust on held-out conditions (expanding the RL pipeline as needed for strech goals).

We anticipate the main challenges to be sample efficiency, reward misspecification, instability across various seeds, multi-agent assignment in decentralized modes, and the dominance of existing non-RL methods (primarily rule-based or model-predictive). As a result, we hope to focus on reproducible experiments with clear baselines and evaluations compared to large hyperparameters sweeps. Furthermore, RL is still attractive here since it can learn control policies from direct interaction without an explicit system model, meaning it can generalize across buildings, seasons, or forecasts. Our main goal is producing a strong baseline with meaningful improvement through feature or reward design, and a breakdown of where learned controllers succeed vs fail. Notably, this problem matters since better coordination of distributed batteries and solar generations have direct real-world implications in reducing electricity costs, lower carbon emissions, and easing peak demand on the grid (electrifying the grid!).

The project will be subdivided among the members. Ammaar will work on environment setup and SAC implementation, Erik on PPO implementation and hyperparameter tuning, and Grace will manage the RBC baseline and evaluation pipeline. We will also follow standard group project protocols like meeting biweekly, maintaing code in a shared repo, and keeping a log to track design choices (note: workflow split-up is more subject to change, with this serving as general idea but can adjust as needed).


{\footnotesize
\bibliographystyle{apalike}
\bibliography{references_citylearn}
}

\end{document}
